\chapter{Services, metrics and postmortems}
\label{ch:screens}

Three shorter screens. Each exists because a different person asks a different
question after --- or before --- an incident.

\section{Services: what can break}

\shot{07-services.png}{The services registry. Every incident is attached to one
of these.}

A \emph{service} is a named part of your system: the payments API, the search
cluster, the mobile gateway. Incidents attach to one.

\note{This is why the incidents list can be filtered by service, and why the
metrics page can say ``the payments API had four incidents this month''. Without
a registry you would have a free-text field, everyone would spell things
differently, and no grouping would ever be reliable. A dropdown of known values
is a small constraint that makes every downstream question answerable.}

Each service carries an owning team and a criticality. Only an administrator can
add or change one --- the registry is reference data, and reference data that
anyone can edit stops being reference data.

\section{Metrics: are we getting better or worse?}

\shot{08-metrics.png}{Reliability metrics over a chosen window.}

Two numbers dominate this page, and both are industry standard.

\begin{center}
\begin{tabular}{@{}>{\raggedright\arraybackslash}p{2.2cm}>{\raggedright\arraybackslash}p{4.4cm}p{7.2cm}@{}}
\toprule
\textbf{Metric} & \textbf{Full name} & \textbf{What it tells you} \\
\midrule
MTTA & Mean time to acknowledge & How long before a human picks it up. Measures your \emph{alerting and staffing}. \\
MTTR & Mean time to resolve & How long from report to fixed. Measures your \emph{systems and tooling}. \\
\bottomrule
\end{tabular}
\end{center}

\note{The distinction matters when deciding what to fix. A bad MTTA and a good
MTTR means alerts are not reaching people --- a rota problem. A good MTTA and a
bad MTTR means people respond promptly and then struggle --- a tooling or
architecture problem. One number cannot tell you which, which is why there are
two.}

\subsection{The numbers are precomputed, not counted on the fly}

\where{api/app/Services/Metrics/}
\begin{lstlisting}[language=PHP]
/**
 * Reliability metrics: MTTA, MTTR, throughput and acknowledgement-SLA
 * ...
 * columns `time_to_acknowledge_seconds` and `time_to_resolve_seconds` are
 * ...
 * `percentile_cont` nor `date_trunc` exists in SQLite, and a metrics endpoint
 */
\end{lstlisting}

When an incident is acknowledged, the elapsed seconds are written onto the
incident row there and then. The metrics page averages a stored number rather
than recomputing durations by scanning the timeline.

\note{Storing a derived value is normally poor practice --- it can drift from
the truth. It is right here because the value is \emph{immutable once written}:
the moment of acknowledgement never changes, so the duration never changes
either. The rule of thumb is that caching a derived value is safe exactly when
its inputs can no longer move.}

\warnbox{That comment mentions SQLite, and it is a warning from experience. The
tests run on SQLite for speed; production runs PostgreSQL.
\fpath{percentile\_cont} and \fpath{date\_trunc} exist only in the latter. Code
written naturally against PostgreSQL passes review, fails in the test suite, and
--- worse --- code written to satisfy SQLite can behave differently in
production. Chapter~\ref{ch:bugs-backend} contains a bug of exactly this
family.}

\section{Postmortems: what did we learn?}

\shot{09-postmortems.png}{Postmortems --- the written record produced after an
incident is closed.}

A postmortem is the document written afterwards: what happened, why, what was
done, and what will change so it does not recur.

\subsection{Postmortems are editable --- and that is not a contradiction}

The timeline cannot be edited. A postmortem can. That is deliberate and worth
being precise about.

\begin{center}
\begin{tabular}{@{}>{\raggedright\arraybackslash}p{3.4cm}>{\raggedright\arraybackslash}p{4.2cm}p{6.2cm}@{}}
\toprule
& \textbf{Timeline} & \textbf{Postmortem} \\
\midrule
What it records & What happened & What we understood \\
Can change? & Never & Yes, until published \\
Why & Facts do not change & Understanding improves \\
\bottomrule
\end{tabular}
\end{center}

\where{api/app/Models/Postmortem.php}
\begin{lstlisting}[language=PHP]
public function isPublished(): bool
{
    return $this->status === PostmortemStatus::Published;
}
\end{lstlisting}

A postmortem is a \emph{draft} while being written, and \emph{published} when
finished. Publishing is the point at which it stops being a working document.

\note{This is the same distinction a newspaper makes between a reporter's notes
and a printed article. Conflating them --- letting people quietly revise the
record of events because their understanding changed --- is exactly the failure
mode the append-only timeline exists to prevent.}

\section{Two screens you will meet by accident}

\subsection{Registering an organisation}

\shot{04-register.png}{Creating an organisation. The first account becomes its
owner.}

There is no ``create an account'' for an existing organisation on this screen ---
registering creates an \emph{organisation} and makes you its administrator.
Everyone else is invited from inside.

\warnbox{This is why a production deployment has no demo logins. The seeder
refuses to run outside development, so on a fresh production database the very
first thing you do is register, and that first registration becomes the owner.
Part~III walks through it.}

\subsection{The page that is not there}

\shot{10-not-found.png}{An unknown URL, inside the application shell.}

Notice the navigation is still present. The 404 is rendered \emph{inside} the
layout, so a mistyped link leaves you somewhere you can navigate away from,
rather than at a dead end.

\note{A single-page application has to be told to do this. The server does not
know which paths are real --- it hands \fpath{index.html} to every unknown URL
and lets the router decide. That is the \fpath{try\_files \$uri \$uri/
/index.html} line in the web container's nginx config, and without it a hard
refresh on \fpath{/incidents/42} would return a genuine 404 while clicking
through to the same page worked fine.}

\section{Check your understanding}

\begin{itemize}[leftmargin=1.4cm]
\cbox Why is a service registry better than a free-text field?
\cbox What does a good MTTA with a bad MTTR suggest you should fix?
\cbox Why is it acceptable to store \fpath{time\_to\_resolve\_seconds} on the
  incident row?
\cbox Why may a postmortem be edited when the timeline may not?
\end{itemize}
